\section*{Chapter \ref{ch:g9:html-and-css} --- Web Pages: HTML and CSS}

\begin{solution}{exo:g9:html-and-css:1}
\code{ul}, then \code{li} indented once, \code{a} indented twice (inside the
first \omterm{def:g7:lists:list}{item}), and the second \code{li} indented once. At the end, nothing is
left open.
\end{solution}

\begin{solution}{exo:g9:html-and-css:2}
\code{body} applies to the body, and so to the whole page shown; \code{h1} to
the main heading; \code{li.new} to the second \omterm{def:g7:lists:list}{item} only, Chess club;
\code{p} to the two paragraphs.
\end{solution}

\begin{solution}{exo:g9:html-and-css:3}
No padding and no border: the box is 300 \omterm{def:g5:text-pictures-sound-as-numbers:pixel}{pixels} wide, and
$300 + 2 \times 10 = 320$ with its margins.
\end{solution}

\begin{solution}{exo:g9:html-and-css:4}
\omcode{code/grade-9/01-html-and-css/exercises/my_page.html}{1}{14}{A page of
one's own.}\label{lst:g9:html-and-css:sol-page}

A test of the chapter checks that its tags pair up.
\end{solution}

\begin{solution}{exo:g9:html-and-css:5}
\code{em} is inside \code{p}, so it must close first:
\code{<p>An <em>important</em> day</p>}, with the tree \code{p}, then
\code{em} below it. \code{show_tree} does not notice the mistake: it only
counts the open \omterm{def:g9:html-and-css:element}{elements}, and two \omterm{def:g9:html-and-css:element}{start tags} followed by two end tags bring
the count back to 0, whatever their names. \Cref{exo:g9:html-and-css:11}
looks at the names.
\end{solution}

\begin{solution}{exo:g9:html-and-css:6}
The \code{alt} text is missing. People who listen to the page with a screen
reader hear nothing about the map, and everyone sees an empty frame when the
image does not load. The fix describes it:
\code{<img src="map.png" alt="Map of the town centre">}.
\end{solution}

\begin{solution}{exo:g9:html-and-css:7}
\omcode{code/grade-9/01-html-and-css/exercises/places.html}{1}{18}{Four
places.}\label{lst:g9:html-and-css:sol-places}
\end{solution}

\begin{solution}{exo:g9:html-and-css:8}
\omcode{code/grade-9/01-html-and-css/exercises/h2.css}{1}{3}{Every
\code{h2} in royal blue.}\label{lst:g9:html-and-css:sol-h2}
\end{solution}

\begin{solution}{exo:g9:html-and-css:9}
10 \omterm{def:g9:html-and-css:element}{elements}: \code{table}, three \code{tr}, two \code{th} and four \code{td}.
Depth 3: \code{table}, then \code{tr}, then the cells.
\end{solution}

\begin{solution}{exo:g9:html-and-css:10}
Keep the \omterm{def:g9:html-and-css:element}{start tags} named \code{a}, and cut out the text between
\code{href="} and the next quote:

\omcode{code/grade-9/01-html-and-css/exercises/links.py}{8}{15}{The addresses of
the \omterm{def:g5:networks-first-look:web}{links}.}\label{lst:g9:html-and-css:sol-links}

On the library page it returns \code{['books.html']}.
\end{solution}

\begin{solution}{exo:g9:html-and-css:11}
Suppose an end tag \code{</a>} comes while an \omterm{def:g9:html-and-css:element}{element} \code{b}, opened after
\code{a}, is still open. Then \code{b} started inside \code{a}, so all of its
content, end tag included, belongs inside \code{a}, before \code{</a>}. But
\code{b} is still open at \code{</a>}: its end comes after, outside \code{a}.
That is impossible in a page whose \omterm{def:g9:html-and-css:element}{elements} sit inside one another; so every
end tag closes the \omterm{def:g9:html-and-css:element}{element} opened most recently and still open, exactly as a
closing bracket closes the last open bracket. The page of
\cref{exo:g9:html-and-css:5} breaks the rule: \code{</p>} comes while
\code{em} is still open. To catch it, a \omterm{def:g4:programs-and-operating-system:program}{program} would keep the \omterm{def:g7:lists:list}{list} of the
names still open, check that each end tag has the name at the end of that
\omterm{def:g7:lists:list}{list}, then remove it.
\end{solution}

\begin{solution}{exo:g9:html-and-css:12}
\omcode{code/grade-9/01-html-and-css/exercises/csv_to_html.py}{1}{17}{A CSV table
turned into \omterm{def:g9:html-and-css:html}{HTML}.}\label{lst:g9:html-and-css:sol-table}

A test turns the town's \code{places.csv} into a table and checks that its
tags pair up and that it has one row per line of the \omterm{def:g3:files-and-folders:file}{file}.
\end{solution}

\begin{solution}{pb:g9:html-and-css:1}
\textbf{1.} \code{html} holds \code{head} (with \code{title}) and
\code{body}; \code{body} holds \code{h1}, \code{p} (with \code{strong}) and
\code{table}; \code{table} holds two \code{tr}; the first holds two
\code{th}, the second two \code{td}.

\textbf{2.} 14 \omterm{def:g9:html-and-css:element}{elements}.

\textbf{3.} Depth 5: \code{html}, \code{body}, \code{table}, \code{tr}, then
the cells. The two \code{th} and the two \code{td} are the deepest.

\textbf{4.} The \omterm{def:g9:html-and-css:element}{element} being opened sits one level below the \omterm{def:g9:html-and-css:element}{elements} already
open: its level is \code{depth + 1}.

\textbf{5.} An \code{img} is an \omterm{def:g9:html-and-css:element}{element} of the tree, but has no end tag:
raising \code{depth} for it would leave the count one too high for the rest of
the page.

\textbf{6.} Yes: 14 \omterm{def:g9:html-and-css:element}{elements} and depth 5.

\omoutput{code/grade-9/01-html-and-css/outputs/page_stats.txt}{Three pages, then
\omterm{def:g7:lists:list}{lists} of 1 to 10 \omterm{def:g7:lists:list}{items}.}\label{out:g9:html-and-css:stats}

\textbf{7.} The text of a page with \code{html}, \code{head}, \code{title},
\code{body}, \code{h1}, \code{ul} and three \code{li}, all on one line: 9
\omterm{def:g9:html-and-css:element}{elements}, the three \omterm{def:g7:lists:list}{items} side by side under \code{ul}.

\textbf{8.} The count goes up by one with each \omterm{def:g7:lists:list}{item}, from 7 to 16; the depth
stays 4.

\textbf{9.} $n + 6$: the six \omterm{def:g9:html-and-css:element}{elements} \code{html}, \code{head}, \code{title},
\code{body}, \code{h1} and \code{ul}, plus one \code{li} per \omterm{def:g7:lists:list}{item}.

\textbf{10.} The \omterm{def:g7:lists:list}{items} sit side by side in the same \code{ul}: a longer \omterm{def:g7:lists:list}{list}
makes the tree wider, not deeper.

\textbf{11.} \textbf{Result.} A page with a \omterm{def:g7:lists:list}{list} of $n$ \omterm{def:g7:lists:list}{items} has $n + 6$
\omterm{def:g9:html-and-css:element}{elements}, and its tree has depth 4 whatever $n$: the six \omterm{def:g9:html-and-css:element}{elements} of the
skeleton are fixed, and each \omterm{def:g7:lists:list}{item} adds one \omterm{def:g9:html-and-css:element}{element} beside the others, never
below them.
\end{solution}
